\documentclass[11pt]{article}

\newcommand{\cnum}{Math 132H}
\newcommand{\ced}{Winter 2026}
\newcommand{\ctitle}[4]{\title{\vspace{-0.5in}\cnum, \ced\\Week #1: #2}\author{\vspace{-0.35in}\\#3, #4}}
\usepackage{enumitem}
\usepackage[usenames,dvipsnames,svgnames,table,hyperref]{xcolor}
\usepackage{amsmath, amsfonts}
\usepackage{graphicx} % Required for inserting images
\usepackage{float}
\usepackage{listings}
\usepackage{xcolor}
\usepackage{hyperref}
\usepackage{comment}

\renewcommand*{\theenumi}{\alph{enumi}}
\renewcommand*\labelenumi{(\theenumi)}
\renewcommand*{\theenumii}{\roman{enumii}}
\renewcommand*\labelenumii{\theenumii.}

\definecolor{codegreen}{rgb}{0,0.6,0}
\definecolor{codegray}{rgb}{0.5,0.5,0.5}
\definecolor{codepurple}{rgb}{0.58,0,0.82}
\definecolor{backcolour}{rgb}{0.95,0.95,0.92}
\definecolor{header}{HTML}{205459}
\definecolor{solution}{HTML}{204d52}

\newcommand{\solution}[1]{{{\textcolor{header}{Solution:} \textcolor{solution}{#1}}}}
\setlist[enumerate]{label=(\alph*)}

\lstdefinestyle{mystyle}{
    backgroundcolor=\color{backcolour},
    commentstyle=\color{codegreen},
    keywordstyle=\color{magenta},
    numberstyle=\tiny\color{codegray},
    stringstyle=\color{codepurple},
    basicstyle=\ttfamily\footnotesize,
    breakatwhitespace=false,
    breaklines=true,
    captionpos=b,
    keepspaces=true,
    numbers=left,
    numbersep=5pt,
    showspaces=false,
    showstringspaces=false,
    showtabs=false,
    tabsize=2
}


\begin{document}
\ctitle{\#4}{}{spongebob squarepants}{krusty krab}
\date{}
\maketitle

Let $\sum_{n=0}^{\infty}\alpha_nz^{n}$ be a powerseries then 
let
\[
\frac{1}{R} = \limsup_{n\to \infty}|\alpha_n|^{\frac{1}{n}}
\] 
power series converges uniformly for $z \le r$ for any  $r < R$ and diverges for  $|z| > R$.
Also
 \[
f^{(k)}(z) = \sum_{n=k}^{\infty}n!\beta_n z^{n-k}
\] 
with
\[
\beta_n = \frac{f^{(n)}(0)}{n!}
\] 
$z : [a,b] \to \mathbb{C}$ and $\tilde z : [c,d] \to \mathbb{C}$ are equivalent if there exists
\[
    t(s) : [c,d] \to^{\text{onto}} [a,b]
\] 
such that $t'(s) > 0$ and  $z(t(s)) = \tilde z(s)$
and
 \[
\int_{\gamma}^{}f(z)dz = \int_{a}^{b}f(z(t))z'(t)dz
\] 
\[
\int_{\gamma}^{}F'(z)dz = F(z(b)) - F(z(a))
\] 

\[
\frac{1}{2\pi i} \int_{\partial D}^{} \frac{f(w)}{w-z}dz = f(z)
\] 
if $z \in D$ and  $f$ is holo in a domain containing the closure of $D$.
also
\[
f^{(n)}(z) = \frac{n!}{2\pi i} \int_{ \partial D}^{}\frac{f(\zeta)}{(\zeta - z)^{n+1}}d \zeta
\] 
\[
    a^{n} - b^{n} = (a-b)(a^{n-1} + a^{n-2}b +\dots + ab^{n-2} + b^{n-1})
\] 
\[
    |f^{(n)}(a)| \le \frac{n!}{r^{n}}\sup_{|z-a|= r}|f(z)|
\] 
\[
u : \Omega \to \mathbb{C}
\] 
is Harmonig if $D = \{|z-z_0| < R\} \subset \overline{ D} \subset \Omega$ and
\[
u(z_0) = \frac{1}{2\pi} \int_{0}^{2\pi} u(z_0 + Re^{ i \theta}) d \theta
\] 
If $f_n$s are holo and $f_n \to f$ uniformly on all compact  $K \subset \Omega$ then $f$ is holomorphic and $f_n^{(k)} \to f^{(k)}$ as well\\
If $F(z,t) $ defined on $\Omega \times [a,b]$ and  $F(z,t)$ continuous on  $\Omega \times [a,b$
and  $F(z,t)$ is holomorphic holding  $t$ constant then 
\[
F(z) = \int_{a}^{b}F(z,t)dt
\] 
is holomorphic for all $z \in \Omega$.
 $f$ holo on upper half plane the  $\overline{f(\overline{z}}$ holo on lower half plane. $f$ holo on the inside of the disc somewhere then $\overline{f(\frac{1}{\overline{ z}}}$ holo where $w = \frac{1}{ \overline{ z}} $ for some $z$ in the domain.

If  $K \subset \mathbb{C}$ is compact, then polynomials in $z$ and $\overline{ z}$ are dense in $ C(K)$ the set of continuous functions from  $K $ to  $ \mathbb{C}$\\
If $f$ is holomorphic on $A = \{ R_0 < |z| < R_1\}$  and $r_0 < |z| < r_1$ then
\[
f(z) = \frac{1}{2\pi i} \int_{|z| = r_1}^{} \frac{f(w)}{w-z}dw - \frac{1}{2\pi i}\int_{|z| = r_0}^{}\frac{f(w)}{w-z}dw
\] 

Runge's theorem: Let $K \subset \mathbb{C}$ be compact and let $f(z)$ be a function holomorphic on  $\Omega \supset K$ open
then therer exists a sequence  $R_n$ of rational functions poles not in  $K$ such that
\[
\sup_K | R_n(z) - f(z)| \to 0
\] 
READ THIS PROFO AGAIN RUNGE's THEOREM
If $K \subset C$ is compact and its complement connected and $f$ is holo on an open set containing  $K$ then you can uniformly approximate $f$ by polynomials on $K$.

\[
\frac{1}{2\pi i} \int_{ \gamma}^{} \frac{f'(z)}{f(z)}dz = N \in \mathbb{Z}
\] 
and the general winding number is when $f(z) = (z-a)$ but
\[
\frac{1}{2\pi i} \int_{\gamma}^{} \frac{f'(z))}{f(z)} = Z - P
\] 
equals number of zeros (counting multiplicities) minus number of poles (also with multiplicities)
The following are equivalent
\begin{enumerate}
    \item $\Omega$ is simply connected
    \item For all holomorphic  $f$ there is a $F$ such that  $F' = f$
    \item For all holo  $f$ on  $\Omega$ such that  $f(z) \ne 0$ there exists  $q$ such that  $f = e^{q}$ on $\Omega$
    \item for all piecewise smooth curves and all  $f$ holomorphic  
         \[
        \int_{ \gamma}^{}f(z_0dz = 0
        \] 
    \item for all piecewise smooth closed curvves $\gamma$ and all $a \not\in \Omega$ 
        \[
        \int_{\gamma}^{} \frac{dz}{z-a} =0
        \] 
\end{enumerate}

$f$ holomorphic on $\{0 < |z-z_0| < R\}$ then
\[
    f(z) = \sum_{n=0}^{\infty}a_n(z-z_0)^{n} - \sum_{n=1}^{\infty}a_{-n}(z-z_0)^{-n}
\] 
with
\[
a_n = \frac{1}{2\pi i} \int_{|\zeta - z_0| = r}^{}\frac{f(\zeta)}{(\zeta - z_0)^{n+1}} d\zeta
\] 
and
\[
    a_{-k} = \frac{1}{2\pi i } \int_{|\zeta - z_0| = r_1}^{} f(\zeta) (\zeta - z_0)^{k} d\zeta
\] 
If $f(z)$ has an essential singularity at  $z_0$ then for every $\epsilon$ $f^{-1}(|z-z_0| < \delta\}$ is dense in $ \mathbb{C}$
\\
If $f(z) = \frac{a_{-N}}{(z-z_0)^{N}} + \dots + a_0 + a_1(z-z_0) + \dots$ then
\[
    a_{-1} = \frac{1}{2\pi i} \int_{ |\zeta - z_0| = r}^{}f(\zeta) d\zeta
\] 
and if $N = 1$ then  $a_{-1} = \lim_{z\to z_0}f(z)$ but in general
\[
    a_{-k} = \frac{1}{(k+N)!}\lim_{z\to z_0} \frac{d^{k+m}}{dz^{k+m}}[(z-z_0)^{m}f(z)]
\] 
$f$ is meromorphic if its set of poles has no limit points in  $\Omega$ and each pole is not essential.
.If $\infty \in \Omega$ then $f$ is meromorphic if $f(\frac{1}{w})$ is meromorphic on $\{|w| < \frac{1}{R}\}$ where $\{|z| > R\} \subset \Omega$\\
If $f$ is meromorphic on $ \mathbb{C}^{*}$ then $f$ is rational.

holomorphic functions are open mappings if they are non constant.

Let $f$ and $g$ be holomorphic functions inside a closed curve $C$. If $|g(z)| < |f(z)|$ for all $z \in C$ then  $f(z)$ and  $f(z) + g(z)$ have the 
same number of zeros inside  $C$ counting multiplicities



\end{document}
